\documentclass[12pt,a4paper]{ctexart}
\usepackage[margin=2.2cm]{geometry}
\usepackage{graphicx}
\usepackage{booktabs}
\usepackage{amsmath}
\usepackage{float}
\usepackage{hyperref}
\usepackage{xcolor}
\usepackage{listings}

\lstset{
  language=Python,
  basicstyle=\ttfamily\small,
  keywordstyle=\color{blue!70!black},
  commentstyle=\color{green!40!black},
  stringstyle=\color{orange!60!black},
  showstringspaces=false,
  breaklines=true,
  frame=single,
  tabsize=4
}

% ===== 模板配置 - 请根据实验要求修改 =====
\newcommand{\reporttitle}{实验报告标题（如：CIFAR-10 分类与预处理）}
\newcommand{\subtask}{示例：两个任务演示}
\newcommand{\authname}{姓名}
\newcommand{\authid}{学号}
% ==========================================

\title{\reporttitle}
\author{\authname：\underline{\hspace{3cm}} \quad \authid：\underline{\hspace{3cm}}}
\date{\today}

\begin{document}
\maketitle

\section{实验目的}
本实验的目标如下：
\begin{itemize}
  \item 目标一：\textbf{【请填入实验第一个目标】}
  \item 目标二：\textbf{【请填入实验第二个目标】}
  \item 目标三：\textbf{【请填入实验第三个目标】}
\end{itemize}

\section{实验过程}
\subsection{实验环境}
\begin{itemize}
  \item 操作系统：\textbf{【如 Linux、Windows、MacOS】}
  \item 编程语言：Python 3
  \item 深度学习框架：\textbf{【如 PyTorch、TensorFlow】}
  \item 其他关键工具：\textbf{【如 NumPy、Matplotlib、scikit-learn】}
\end{itemize}

\subsection{数据集与预处理}
\begin{itemize}
  \item \textbf{【填入数据集名称与规模】}
  \item 数据特征：\textbf{【输入维度、输出类别、数据格式等】}
  \item 预处理步骤：\textbf{【标准化、增强、分割等】}
\end{itemize}

\subsection{实现步骤}
\begin{enumerate}
  \item \textbf{【步骤一描述】}
  \item \textbf{【步骤二描述】}
  \item \textbf{【步骤三描述】}
\end{enumerate}

\section{实验内容}
\subsection{任务一：【任务一的名称】}
\textbf{【对任务的详细描述，包括方法、参数、公式等】}

主要参数：
\begin{itemize}
  \item 参数一：\textbf{【取值与说明】}
  \item 参数二：\textbf{【取值与说明】}
\end{itemize}

实现要点：
\begin{itemize}
  \item 关键实现点一
  \item 关键实现点二
\end{itemize}

\subsection{任务二：【任务二的名称】}
\textbf{【对任务的详细描述】}

\textbf{【如有公式，示例如下：】}
\begin{equation}
y = f(x)
\end{equation}

\section{实验结果}
\subsection{任务一结果：【结果名称】}
\begin{figure}[H]
\centering
\includegraphics[width=0.95\textwidth]{images/【图像文件名】}
\caption{【图像标题】}
\label{fig:task1}
\end{figure}

结果分析：
\begin{itemize}
  \item 结果特性一
  \item 结果特性二
  \item 结果特性三
\end{itemize}

\subsection{任务二结果：【结果名称】}
\begin{table}[H]
\centering
\caption{【表格标题】}
\label{tab:task2}
\begin{tabular}{lccc}
\toprule
指标 & 值1 & 值2 & 值3 \\
\midrule
指标A & 值 & 值 & 值 \\
指标B & 值 & 值 & 值 \\
\bottomrule
\end{tabular}
\end{table}

结果分析：
\begin{itemize}
  \item 分析点一
  \item 分析点二
\end{itemize}

\section{总结}
本实验成功完成了【主要任务总结】：
\begin{itemize}
  \item 完成内容一
  \item 完成内容二
  \item 完成内容三
\end{itemize}

\textbf{【可选：经验或建议】}

\section*{附录：实验源代码}
\subsection*{A.1 主要源码：【源文件名.py】}
\begin{lstlisting}
【粘贴完整源代码】
\end{lstlisting}

\subsection*{A.2 运行命令}
\begin{verbatim}
# 任务一运行命令
python script1.py --arg1 value1 --arg2 value2

# 任务二运行命令
python script2.py --arg1 value1
\end{verbatim}

\end{document}

% ===== 使用说明 =====
% 1. 修改顶部的 \newcommand 部分填入实验基本信息
% 2. 替换所有【】括号内的占位符为实际内容
% 3. 在 images/ 目录下放置实验结果图像
% 4. 在附录中粘贴完整源代码
% 5. 使用 xelatex report.tex 进行编译
% ====================
